% !TEX root = ../main.tex
\chapter{Hardening and Edge Cases}
\label{ch:compile_hardening}

The 3.5 compile path went through a substantial hardening pass
between the initial 3.5-rc1 release and the final 3.5.0 release.
This chapter catalogues the edge cases that were stabilised:
the dynamic-shape handling (deferred to 3.5.1), the
unsupported-op fallback, the NaN propagation under compile, the
MPSGraph SDK version sensitivity, and the test infrastructure
that catches regressions.

\section{Dynamic Shapes: The Deferred Goal}
\label{sec:compile_hardening:dyn_shapes}

The largest goal that did not make 3.5: dynamic-shape support.
A compiled graph that accepts variable input shapes (e.g., a
batch size that changes per call without recompilation) would
let workloads with naturally variable shapes benefit from
compilation.

\subsection{Why deferred}
\label{ssec:compile_hardening:dyn_deferred}

\MPSGraph\ provides a dynamic-shape mode (placeholders with
symbolic dimensions), but the integration with
\code{gradientForPrimaryTensor} produced shape mismatches in
representative models. The investigation found:

\begin{itemize}
  \item Several backward formulas rely on concrete shapes for
    correct broadcasting (e.g., the unbroadcast step requires
    knowing which dimensions were size 1).
  \item Some emitters use \code{shape}-derived attributes
    (e.g., \code{layer\_norm} normalised\_shape) that cannot be
    symbolic.
  \item The parity tests on GPT-2-base and ViT-B failed under
    symbolic primary shapes.
\end{itemize}

The fix is per-op work (each emitter needs review for
symbolic-shape correctness) and is targeted for 3.5.1. The 3.5
release ships shape-specialised compilation only.

\subsection{User-side workaround}
\label{ssec:compile_hardening:dyn_workaround}

For shape-variable workloads, the recommended pattern is to pad
to a fixed shape. The padding adds wasted compute but is
typically less than the recompile cost.

For language modelling with variable sequence lengths: pad to
the maximum length and use attention masking. For batched
inference with variable batch sizes: dispatch to per-size
caches (the user maintains multiple compiled modules, one per
expected batch size).

\section{Unsupported-Op Fallback}
\label{sec:compile_hardening:unsupported_fallback}

The compile path falls back to eager when an unsupported op
appears in the trace
(\chref{ch:compile_design}~Section~\ref{ssec:compile_design:goal_fallback}).
The hardening work focused on making the fallback robust.

\subsection{The connected-components analysis}
\label{ssec:compile_hardening:cc_analysis}

The MpsBuilder partitions the trace into compiled and eager
components via a connected-components analysis on the
unsupported-op DAG. Each unsupported op is a fragmentation
point; the surrounding ops are clustered into the maximum
compiled regions.

\subsection{Partial fallback bugs}
\label{ssec:compile_hardening:partial_fallback}

The 3.5-rc1 implementation had several bugs in the partial-
fallback path: tensor lifetime management between compiled and
eager regions, dtype mismatches at the boundary, gradient
propagation through the fallback boundary. The 3.5.0 release
fixed all known cases; the regression test suite verifies
correctness on roughly 30 model configurations involving
unsupported ops.

\subsection{The strict-mode safety net}
\label{ssec:compile_hardening:strict_safety}

\code{compile(model, strict=True)} disables the fallback. For
model authors writing new models, the recommendation is to
develop under \code{strict=True} so that unsupported ops surface
immediately; switch to \code{strict=False} for production
training.

\section{NaN Handling Under Compile}
\label{sec:compile_hardening:nan_handling}

In eager mode, a NaN gradient can be caught by
\code{set\_detect\_anomaly(True)}, which records the
forward-call stack and reports the offending op. Under compile,
the per-op stack is lost (the compiled graph executes as one
unit), so anomaly detection is less informative.

\subsection{The mitigation}
\label{ssec:compile_hardening:nan_mitigation}

The 3.5 release added a \code{check\_nan} flag:

\begin{lstlisting}[style=lucid,language=LucidPython]
compiled = lucid.compile(model, check_nan=True)
\end{lstlisting}

When enabled, the compiled graph includes per-output NaN checks.
The cost is a few percent overhead; the benefit is that a NaN
surfaces at the compiled-call boundary with a diagnostic message
naming which output has the NaN.

For deeper diagnosis (which sub-op produced the NaN), the user
re-runs in eager with the anomaly detector.

\section{MPSGraph SDK Version Sensitivity}
\label{sec:compile_hardening:sdk_sensitivity}

\MPSGraph's API surface changes per macOS release. The 3.5
release targets macOS 26.0 \MPSGraph, which exposes the SDPA
fused kernel and the gradient-construction method.

\subsection{Per-op feature gates}
\label{ssec:compile_hardening:feature_gates}

Some emitters use SDK-conditional code paths (e.g., the SDPA
emitter falls back to the composite path if the fused kernel is
unavailable). The emitter dispatch is at module load time, so
\Lucid\ pins to the minimum-target SDK and assumes the features
are available.

Users on a non-supported macOS see a clear error at engine
load, not at runtime.

\subsection{macOS version requirement}
\label{ssec:compile_hardening:macos_req}

\Lucid\ 3.5 requires macOS 26.0 or newer. Older macOS versions
lack the \MPSGraph\ features that the compile path uses. The
\code{pyproject.toml} declares this:

\begin{lstlisting}[style=lucid,language=LucidPython]
[project]
requires-python = ">=3.14"
classifiers = [
    "Operating System :: MacOS",
]
\end{lstlisting}

Plus the wheel tag \code{macosx\_26\_0\_arm64} (matching
\MLX's) prevents installation on older macOS.

\section{The Compile-Time Cost Profile}
\label{sec:compile_hardening:compile_cost}

The hardening work also reduced compile-time cost. Comparing
3.5-rc1 to 3.5.0 final:

\begin{itemize}
  \item GPT-2-base initial compile: 4.8 s $\to$ 1.8 s.
  \item ViT-B initial compile: 2.1 s $\to$ 0.9 s.
  \item ResNet-50 initial compile: 1.4 s $\to$ 0.7 s.
\end{itemize}

The reduction came from emitter optimisations (avoid redundant
intermediate \MPSGraph\ nodes), trace-IR simplifications (skip
nodes that are no-ops in the compile target), and parallelising
the per-op shape inference.

\section{The Regression Test Suite}
\label{sec:compile_hardening:regression_tests}

The compile path's regression tests live in
\code{lucid/test/compile/}. They cover:

\begin{itemize}
  \item Per-op parity: each emitter is tested on a one-op model
    with diverse shapes and dtypes.
  \item Per-layer parity: each standard layer (Linear, Conv,
    LayerNorm, MHA) is tested as a compiled module.
  \item Per-model parity: each model-zoo entry runs forward +
    backward both compiled and eager, asserting agreement to
    $10^{-5}$ relative tolerance.
  \item Fallback correctness: models with deliberately
    unsupported ops are tested under \code{strict=False} to
    verify the partial-fallback path produces correct results.
  \item Cache behaviour: multi-shape workloads test cache hits,
    eviction, and recompile.
\end{itemize}

The test suite runs in CI on every PR; a parity failure blocks
the merge.

\section{Known Limitations (3.5)}
\label{sec:compile_hardening:known_limits}

We document the limitations explicitly so users do not waste
time discovering them.

\begin{itemize}
  \item Dynamic shapes: not supported. Pad to fixed shape.
  \item Custom Functions (\code{lucid.autograd.Function} subclasses):
    not compiled; trigger eager fallback for their region.
  \item Linalg decompositions (\code{cholesky}, \code{qr}, etc.):
    no emitter; trigger eager fallback.
  \item FFT: no emitter (will be added in a future release).
  \item Sparse tensors: not supported anywhere in \Lucid\ 3.5.
  \item Control flow with tensor-dependent branches: traced
    incorrectly; user must rewrite using \code{where}.
  \item Data-dependent output shape (\code{nonzero},
    \code{unique}, \code{masked\_select}): always fall back to
    eager.
\end{itemize}

\section{Future Work}
\label{sec:compile_hardening:future_work}

The 3.5.1 release plan:

\begin{enumerate}
  \item Dynamic-shape support for the common cases
    (batch-variable, sequence-variable).
  \item FFT emitters.
  \item Whole-program compile: fuse across the optimiser step.
  \item Improved fallback granularity: per-op rather than
    per-region partition.
\end{enumerate}

Beyond 3.5.1: SME-based direct emission
(\chref{ch:apple_silicon}~Section~\ref{ssec:apple_silicon:sme})
that bypasses \MPSGraph\ for some operations. This would let
\Lucid\ optimise for hardware features not yet exposed by Apple's
graph library.

\section{Summary and Forward References}
\label{sec:compile_hardening:summary}

The 3.5 compile path is production-stable for shape-fixed
training workloads. Unsupported ops fall back gracefully; NaN
detection is available via \code{check\_nan}; the regression
test suite enforces parity across model configurations. Known
limitations (dynamic shapes, custom Functions, linalg, FFT)
are explicit and documented. The 3.5.1 release will close
several gaps.

This chapter closes \partref{part:graph_compile}.
\partref{part:model_zoo} treats the model zoo --- the catalogue
of ready-to-use architectures that benefit from the compile
path.

\begin{lucidnote}
For the user: the compile path is opt-in and reversible. If a
model trains correctly in eager but fails under compile, drop
\code{lucid.compile()} and report the issue. The fallback
mechanism is designed so users rarely encounter compile-only
failures in practice.
\end{lucidnote}
